\documentclass{article}
\usepackage{../../math_template}

\author{Jonathan Kasongo}
\title{Noether Sheet 5}

\begin{document}
\maketitle
\centerline{
These were my solutions to the Noether Mentoring Scheme Sheet 5. Problem
statements are abridged.
}

\begin{problem}{Problem 1}
(a) Suppose that every positive integer is coloured red or blue. Show that
there must be integers $u,v, w$ and $x$ (which need not all be different)
of the same colour such that $u + v + w = x$. (This is called a
\textit{monochromatic} sum.)
\vspace{0.2cm}

(b) What is the largest $N$ such that the integers $1, \ldots, N$ can be
coloured red and blue without the sum $u+v+w=x$ having a monochromatic
solution?
\end{problem}

\begin{solution}{Solution}
(a) Let us assume WLOG that $1$ is blue. We can do this because if $1$ was red
then the arguement follows the same we just exchange the words red and blue.
We argue by contradiction.
\vspace{0.2cm}

Note that if $k$ and $3k$ are the same colour than we have a contradiction
because $(k) + (k) + (k) = (3k)$ is a monochromatic sum.
\vspace{0.2cm}

Suppose that we have found a blue integer $k > 2$ such that $k + 1$ is red.
Since $k$ is blue, $3k$ must be red. Because there doesn't exist a
monochromatic sum there cannot exist a red integer $y$ such that
$$
(k+1) + (k+1) + (y) = (3k) \iff y = k - 2
$$
and so $k-2$ must be blue. But now observe that
$$
(k-2) + (1) + (1) = (k)
$$
is a monochromatic sum where every integer is blue, contradiction.
\vspace{0.2cm}

If we cannot find such a blue integer $k$ then either the numbers
$3,4,5,\ldots$ must be coloured
$$
\text{red, red,} \ldots \text{red, blue, blue, blue,} \ldots
$$
or such that every one of the numbers $3,4,5\ldots$ is blue, or such that
every one of those numbers is red. These are the only constructions,
because any other construction would have a point at which $k$ is blue
and $k+1$ is red.
In each case the colour
of the numbers is eventually constant so clearly for large enough $k$
we can find integers $k, \ 3k$ that are the same colour, then we are done
because of the reasoning presented in the second paragraph. $\square$
\\

(b) Again we can assume that $1$ is blue. We look at two cases.
\vspace{0.2cm}

\textit{Case 1:} If $2$ is blue, then we have
\begin{center}
\begin{tblr}{cc}
1 & 2 & 3 & 4 & 5 & 6 & 7 & 8 & 9 & 10 & 11 & 12\\ \hline
B & B \\
\end{tblr}
\end{center}
Now if $k$ is blue you can make the monochromatic sums
$(k) + (1) + (1), \ (k) + (1) + (2), \ (k) + (2) + (2)$. So if $k$ is blue
then $k+2, k+3$ and $k+4$ must all be red. This gives
\begin{center}
\begin{tblr}{cc}
1 & 2 & 3 & 4 & 5 & 6 & 7 & 8 & 9 & 10 & 11 & 12 \\ \hline
B & B & R & R & R & R & \\
\end{tblr}
\end{center}
Now note that $k$ and $3k$ must be different colours. Also since
$3$ is red, if $k$ is red then you can make the monochromatic sum
$(k) + (3) + (3)$. So if $k$ is red then $k+6$ must be blue. This gives
\begin{center}
\begin{tblr}{cc}
1 & 2 & 3 & 4 & 5 & 6 & 7 & 8 & 9 & 10 & 11 & 12\\ \hline
B & B & R & R & R & R &   &   & B & B  & B  & B \\
\end{tblr}
\end{center}
But now since $(9) + (1) + (1) = (11)$ is a monochromatic sum with all
numbers being blue, and we previously deduced that the number $11$ must
be blue, we have a contradiction. So in this case
$N \leq 10$. In particular $N = 10$ is achievable with construction
\begin{center}
\begin{tblr}{cc}
1 & 2 & 3 & 4 & 5 & 6 & 7 & 8 & 9 & 10 \\ \hline
B & B & R & R & R & R & R & R & B & B  \\
\end{tblr}
\end{center}
which can be checked to work.
\vspace{0.2cm}

\textit{Case 2:} If $2$ is red, then we have
\begin{center}
\begin{tblr}{cc}
1 & 2 & 3 & 4 & 5 & 6 & 7 & 8 & 9\\ \hline
B & R
\end{tblr}
\end{center}

Now since $k$ and $3k$ cannot be the same colour, we have

\begin{center}
\begin{tblr}{cc}
1 & 2 & 3 & 4 & 5 & 6 & 7 & 8 & 9\\ \hline
B & R & R &   &   & B
\end{tblr}
\end{center}
But now since $(6) + (1) + (1) = (8)$ is a monochromatic sum with
all numbers being blue, and $(3) + (3) + (2) = (8)$ is a monochromatic sum
with all numbers being red, the number $8$ can't be red or blue. So in
this case $N \leq 7$.\vspace{0.2cm}

Thus we have demonstrated that the largest such $N$ is $N = 10$.
$\blacksquare$

\end{solution}
\newpage

\begin{problem}{Problem 2}
A rectangular piece of paper has a line drawn down its centre.
One corner is folded such that the corner lies on the line, as shown
in the diagram. (Here the top left corner is folded such that
point $A$ is sent to point $M$.)

\begin{center}
\begin{asy}
/* Geogebra to Asymptote conversion, documentation at artofproblemsolving.com/Wiki go to User:Azjps/geogebra */
import graph; size(170);
real labelscalefactor = 0.5; /* changes label-to-point distance */
pen dps = linewidth(0.7) + fontsize(10); defaultpen(dps); /* default pen style */
pen dotstyle = black; /* point style */
real xmin = -7.319930920336964, xmax = 7.316908081987427, ymin = -4.71504179447276, ymax = 5.1086344198213745;  /* image dimensions */
pen qqzzff = rgb(0.,0.6,1.);
pair A = (-4.,3.), B = (4.,3.), D = (-4.,-3.), C = (4.,-3.), M = (1.1961524227066318,0.);

filldraw((-0.5358983848622455,3.)--M--D--cycle, qqzzff + opacity(0.10000000149011612),  qqzzff);
filldraw((-0.5358983848622455,3.)--A--D--cycle, gray + opacity(0.10000000149011612),  gray);
/* draw figures */
draw(A--B);
draw(B--C);
draw(C--D);
draw(D--A);
draw((-4.,0.)--(4.,0.),  linetype("4 4"));
draw(D--M);
draw((-0.5358983848622455,3.)--M);
draw((-0.5358983848622455,3.)--D);
draw((-0.5358983848622455,3.)--M,  qqzzff);
draw(M--D,  qqzzff);
draw(D--(-0.5358983848622455,3.),  qqzzff);
draw((-0.5358983848622455,3.)--A,  gray);
draw(A--D,  gray);
draw(D--(-0.5358983848622455,3.),  gray);
/* dots and labels */
dot(A,dotstyle);
label("$A$", (-3.9540061460196694,3.113200188792878), NE * labelscalefactor);
dot(B,dotstyle);
label("$B$", (4.038694702440616,3.113200188792878), NE * labelscalefactor);
dot(D,dotstyle);
label("$D$", (-3.9540061460196694,-2.895030352985231), NE * labelscalefactor);
dot(C,dotstyle);
label("$C$", (4.038694702440616,-2.895030352985231), NE * labelscalefactor);
dot(M, dotstyle);
label("$M$", (1.2428939941314627,0.08715706921120271), NE * labelscalefactor);
dot((-0.5358983848622455,3.), dotstyle);
label("$N$", (-0.48940605258558134,3.0912723401002573), NE * labelscalefactor);
clip((xmin,ymin)--(xmin,ymax)--(xmax,ymax)--(xmax,ymin)--cycle);
/* end of picture */
\end{asy}
\end{center}

What is angle $\angle ADN$?
\end{problem}

\begin{solution}{Solution}
We note that since the top left corner is folded, the segment
$AD$ is sent to $DM$, whilst $AN$ is sent to $NM$. This means that
$AD = DM$ and $AN = NM$. Now also since $DN$ is
common to both triangles $\triangle AND$ and $\triangle DMN$, we deduce that
those two triangles are congruent by the SSS criterion. Now draw in
the segment $AM$ like so.

\begin{center}
\begin{asy}
/* Geogebra to Asymptote conversion, documentation at artofproblemsolving.com/Wiki go to User:Azjps/geogebra */
import graph; size(170);
real labelscalefactor = 0.5; /* changes label-to-point distance */
pen dps = linewidth(0.7) + fontsize(10); defaultpen(dps); /* default pen style */
pen dotstyle = black; /* point style */
real xmin = -7.319930920336964, xmax = 7.426547325450532, ymin = -4.71504179447276, ymax = 5.1086344198213745;  /* image dimensions */
pen qqzzff = rgb(0.,0.6,1.); pen ffwwqq = rgb(1.,0.4,0.);
pair A = (-4.,3.), B = (4.,3.), D = (-4.,-3.), C = (4.,-3.), M = (1.1961524227066318,0.);

filldraw((-0.5358983848622455,3.)--M--D--cycle, qqzzff + opacity(0.10000000149011612),  qqzzff);
filldraw((-0.5358983848622455,3.)--A--D--cycle, gray + opacity(0.10000000149011612),  gray);
/* draw figures */
draw(A--B);
draw(B--C);
draw(C--D);
draw(D--A);
draw((-4.,0.)--(4.,0.),  linetype("4 4"));
draw(D--M);
draw((-0.5358983848622455,3.)--M);
draw((-0.5358983848622455,3.)--D);
draw((-0.5358983848622455,3.)--M,  qqzzff);
draw(M--D,  qqzzff);
draw(D--(-0.5358983848622455,3.),  qqzzff);
draw((-0.5358983848622455,3.)--A,  gray);
draw(A--D,  gray);
draw(D--(-0.5358983848622455,3.),  gray);
draw(A--M,  ffwwqq);
/* dots and labels */
dot(A,dotstyle);
label("$A$", (-3.9540061460196694,3.113200188792878), NE * labelscalefactor);
dot(B,dotstyle);
label("$B$", (4.038694702440616,3.113200188792878), NE * labelscalefactor);
dot(D,dotstyle);
label("$D$", (-3.9540061460196694,-2.895030352985231), NE * labelscalefactor);
dot(C,dotstyle);
label("$C$", (4.038694702440616,-2.895030352985231), NE * labelscalefactor);
dot(M, dotstyle);
label("$M$", (1.2428939941314627,0.08715706921120271), NE * labelscalefactor);
dot((-0.5358983848622455,3.), dotstyle);
label("$N$", (-0.48940605258558134,3.0912723401002573), NE * labelscalefactor);
clip((xmin,ymin)--(xmin,ymax)--(xmax,ymax)--(xmax,ymin)--cycle);
/* end of picture */
\end{asy}
\end{center}

Now since $A$ is just the reflection of $D$ over that horizontal
dashed line, and $M$ lies on that dashed line it must be the case that
$AM = DM$. So we can deduce that $\triangle AMD$ is in fact equilateral.
\vspace{0.2cm}

Moreover since $\triangle AND$ is congruent to $\triangle DMN$ we have
$\angle ADN = \angle NDM$. So we deduce that
$$
\angle ADN = \frac12 \angle ADM = \frac12 (60) = 30^\circ. \ \blacksquare
$$
\end{solution}
\newpage

\begin{problem}{Problem 3}
A domino in this question covers exactly two adjacent squares of a standard
$8 \times 8$ chessboard.

\begin{enumerate}[label=(\alph*)]
  \item Can the entire chessboard be covered in dominoes, with every domino
  covering exactly two squares?
  \item One corner is removed from the chessboard. Can it now be covered in
  dominoes?
  \item Two opposite corners are removed. Can it now be covered in dominoes?
  \item If two squares are removed from the chessboard, under what
  condition can the board still be covered in dominoes?
\end{enumerate}
\end{problem}

\begin{solution}{Solution}
(a) The answer is yes. Here is one such way to do so.

\begin{center}
\begin{asy}
usepackage("chessfss"); // optional, for later piece-glyphs if desired

unitsize(0.7cm);   // adjust printed size
real s = 1;        // square side in coordinate units (keep = 1 for simplicity)

// board pens
pen darkSquare = black;
pen lightSquare = white;
pen borderPen = darkSquare + linewidth(0.6);

// circle (domino mark) style
real radius = 0.30; // circle radius (fraction of square side)
pen redFill  = rgb(0.86, 0.20, 0.20) + opacity(0.70);
pen blueFill = rgb(0.15, 0.45, 0.85) + opacity(0.70);
pen circleBorder = black + linewidth(0.6);

// draw the chessboard (a1 = (0,0) dark)
for (int i = 0; i < 8; ++i) {
  for (int j = 0; j < 8; ++j) {
    path square = (i,j)--(i+1,j)--(i+1,j+1)--(i,j+1)--cycle;
    if ((i + j) % 2 == 0) fill(square, darkSquare);
    else fill(square, lightSquare);
    draw(square, borderPen);
  }
}

// draw domino-center circles ON TOP of board
for (int i = 0; i < 8; ++i) {
  for (int j = 0; j < 8; j += 2) {

    int k = (int)(j/2);  // domino row: 0,1,2,3
    pen useFill = ((i + k) % 2 == 0) ? redFill : blueFill;

    pair c1 = (i + 0.5, j + 0.5);
    pair c2 = (i + 0.5, j + 1.5);

    fill(circle(c1, radius), useFill);
    //draw(circle(c1, radius), circleBorder);
    fill(circle(c2, radius), useFill);
    //draw(circle(c2, radius), circleBorder);
  }
}
\end{asy}
\end{center}

(b) The answer is no. We say that two tiles are \textit{neighbours}
if they share a vertical or horizontal border. Each domino covers one
pair of neighbours, that is exactly one white square and exactly one black
square. So if the board can be completely covered in dominoes, then there
has
to be $k$ white tiles and $k$ black tiles. Since one corner is removed
from the chess board there is a different number of white and black tiles
on the board, which means that tiling it with dominoes is now impossible.
$\square$
\\

(c) We assume that the question means two vertically/horizontally opposite
corners. The tiling can't be done if we are speaking about diagonally
opposite corners because of the same reasoning that was presented in (b).
The answer to the first interpretation is yes. It suffices to show one
construction for one pair of vertically opposite corners, because we can
rotate the board to get the other 3 configurations of
vertically/horizontally opposite corners being removed.

\begin{center}
\begin{asy}
usepackage("chessfss"); // optional (keeps file ready for piece glyphs)

unitsize(0.7cm);

// board pens
pen darkSquare = black;
pen lightSquare = white;
pen borderPen = darkSquare + linewidth(0.6);

// circle (domino mark) style
real radius = 0.30;
pen redFill  = rgb(0.86, 0.20, 0.20) + opacity(0.70);
pen blueFill = rgb(0.15, 0.45, 0.85) + opacity(0.70);
pen circleBorder = black + linewidth(0.6);

// helper: does the square (file i, rank j) exist (i=0..7, j=0..7)?
bool squareExists(int i, int j) {
  // remove a1 (i=0,j=0) and a8 (i=0,j=7)
  if (i == 0 && (j == 0 || j == 7)) return false;
  return true;
}

// draw the chessboard, skipping a1 and a8
for (int i = 0; i < 8; ++i) {
  for (int j = 0; j < 8; ++j) {
    if (!squareExists(i,j)) continue;
    path square = (i,j)--(i+1,j)--(i+1,j+1)--(i,j+1)--cycle;
    if ((i + j) % 2 == 0) fill(square, darkSquare);
    else fill(square, lightSquare);
    draw(square, borderPen);
  }
}

// --- A-file: iterate j = 1, 3, 5 (dominoes covering (1,2), (3,4), (5,6)) ---
// A-file coloring: blocks b b r r b b on a2..a7
int i = 0;
for (int j = 1; j <= 5; j += 2) {
  // block index 0,1,2 for j=1,3,5 respectively
  int blockIndex = (int)((j - 1) / 2);
  pen useFill = (blockIndex % 2 == 0) ? blueFill : redFill; // 0->blue,1->red,2->blue

  pair c1 = (i + 0.5, j + 0.5);
  pair c2 = (i + 0.5, j + 1.5);

  if (squareExists(i, j))   { fill(circle(c1, radius), useFill); }
  if (squareExists(i, j+1)) { fill(circle(c2, radius), useFill); }
}

// --- B..H files: iterate i=1..7 and j = 0,2,4,6 (standard domino tiling)
// color by checker on (column i) + (domino-row k = j/2) so horizontally/vertically adjacent differ
for (int i2 = 1; i2 < 8; ++i2) {
  for (int j = 0; j < 8; j += 2) {
    int k = (int)(j/2);  // domino-row index 0..3
    pen useFill = ((i2 + k) % 2 == 0) ? redFill : blueFill;

    pair c1 = (i2 + 0.5, j + 0.5);
    pair c2 = (i2 + 0.5, j + 1.5);

    if (squareExists(i2, j))   { fill(circle(c1, radius), useFill); }
    if (squareExists(i2, j+1)) { fill(circle(c2, radius), useFill); }
  }
}
\end{asy}
\end{center}

(d) Again because of the reasoning in (b), we must remove exactly
1 black square and exactly 1 white square. Now we will prove that this
condition is necessary and sufficient.
\vspace{0.2cm}

It will be useful to us to
introduce some notation to refer to which square we are talking about
on a chess board. Let square $(i,j)$ refer to the square on the $i$-th
row from the top of the board and the $j$-th column from the left of the
board. For example square $(1,1)$ refers to the very top left square.
\vspace{0.2cm}

For some square $(i,j)$ define the \textit{weight} of that square to be
the quantity $i+j \ (\text{mod } 2)$. \\

\textit{Claim 1: Two squares are the same colour, if and only if,
they have the same weight.}
\vspace{0.2cm}

\textit{Proof:} Consider two squares $(i,j)$ and $(x,y)$. We note that
if two squares are neighbours, then they have different colours, and also
the two squares have different weights.
We can
travel from $(i,j) \rightarrow (x,y)$ by moving exactly $|x-i|$ squares
vertically and $|y-j|$ squares horizontally. In total we have travelled
$|x-i| + |y-j|$ squares.
\begin{itemize}
\item If we travel an even number of squares, then $(i,j)$ and $(x,y)$ are
the same colour, because the coloured flipped an even number of times
along our path, and
$|x-i| + |y-j| \equiv (x-i) + (y-j) \equiv 0 \ (\text{mod } 2) \implies
x + y \equiv i + j \ (\text{mod } 2)$. Which means that $(i,j)$ and $(x,y)$
have the same weight.
\item If we travel an odd number of squares, then $(i,j)$ and $(x,y)$
are different colours, since the colours flipped an odd number of times
along our path, and $|x-i| + |y-j| \equiv (x-i) + (y-j) \equiv 1 \
(\text{mod } 2) \implies x+y \nequiv i+j \ (\text{mod } 2)$. Which means
that $(i,j)$ and $(x,y)$ have different weights. $\square$
\end{itemize}
\vspace{0.2cm}

\textit{Claim 2: Let $n, m$ be positive integers, with at least one of
$n, m$ even.
Consider a $n \times m$ chessboard. We can always tile it with dominoes.}
\vspace{0.2cm}

\textit{Proof:} We assume WLOG that $n$ is even. The
construction is simple, for each of the $m$ rows tile the row vertically
with $n/2$ dominoes starting from the very top of the row.
$\square$
\\

\textit{Claim 3: Let $n, m$ be positive integers of different parity.
Consider a $n \times m$ chessboard. If the squares $(1,1)$ and $(n,m)$ are
removed then, it is always possible to tile the resultant board with
dominoes.} \vspace{0.2cm}

\textit{Proof:} We can assume WLOG that $n$ is even and $m$ is odd. Since
$m$ is odd and we've removed the top left square and the bottom right
square, there remains $m-1$ squares on the top-most and bottom-most rows
of the chess board. Since $m-1$ is even we can tile these two rows
horizontally with dominoes. Now we are just left with a $(n-2) \times m$
chessboard to tile, but this can be done due to claim 2. $\square$
\\

Here's an example illustrating claim 3 for the case $(n,m) = (4,3)$.
\begin{center}
\begin{asy}
// 4x3 board (3 columns x 4 rows) with a4 and c1 removed
// Pattern (top->bottom):
// _  b  b
// g  r  g
// g  r  g
// b  b  _

unitsize(0.8cm);

// pens (use pen types; rgb returns a pen)
pen borderPen   = black + linewidth(0.6);
pen redFill     = rgb(0.86,0.20,0.20) + opacity(0.85);
pen blueFill    = rgb(0.15,0.45,0.85) + opacity(0.85);
pen greenFill   = rgb(0.12,0.6,0.25)  + opacity(0.85);
pen squareFill  = white;

// geometry
int cols = 3;   // files a,b,c -> i = 0..2
int rows = 4;   // ranks 1..4     -> j = 0..3

real r = 0.32;  // circle radius (in square units)

// helper: whether square (i,j) exists
bool squareExists(int i, int j) {
  // remove a4 (i=0,j=3) and c1 (i=2,j=0)
  if (i == 0 && j == 3) return false;
  if (i == 2 && j == 0) return false;
  return true;
}

// draw grid squares and borders (white squares)
for (int i = 0; i < cols; ++i) {
  for (int j = 0; j < rows; ++j) {
    if (!squareExists(i,j)) continue;
    path sq = (i,j)--(i+1,j)--(i+1,j+1)--(i,j+1)--cycle;
    if ((i + j) % 2 == 0) squareFill = black;
    else squareFill = white;
    fill(sq, squareFill);
    draw(sq, borderPen);
  }
}

// draw colored circles according to the specified pattern
for (int i = 0; i < cols; ++i) {
  for (int j = 0; j < rows; ++j) {
    if (!squareExists(i,j)) continue;

    pen useFill;

    // Map the pattern (j=3 is top row, j=0 is bottom)
    if (j == 3) {                 // top row: _ b b
      // a4 removed already skipped
      useFill = (i == 1 || i == 2) ? blueFill : blueFill; // i==1 or 2 => blue
    }
    else if (j == 2) {            // second row from top: g r g
      useFill = (i == 1) ? redFill : greenFill;
    }
    else if (j == 1) {            // third row: g r g
      useFill = (i == 1) ? redFill : greenFill;
    }
    else {                        // j == 0 bottom row: b b _
      // c1 removed already skipped
      useFill = (i == 0 || i == 1) ? blueFill : blueFill; // i==0 or 1 => blue
    }

    pair center = (i + 0.5, j + 0.5);
    fill(circle(center, r), useFill);
    //draw(circle(center, r), borderPen);
  }
}
\end{asy}
\end{center}

\textit{Claim 4: Let $a,b$ be postive integers of different parity.
Let $n$ be a positive integer satsifying $\max\{a,b\} \leq 2n$.
Consider a $2n \times 2n$ chessboard with a $a \times b$ sub-chessboard
removed from the board. The resultant shape can be tiled with dominoes.}
\vspace{0.2cm}

\textit{Proof:} WLOG let $a$ be even, and $b$ be odd. Consider squares
on the same row as the $a \times b$ hole in the chessboard. These squares
either fall on the left or on the right of the hole. All of the squares that
fall on one particular side of the hole can be tiled, because these squares
will make a sub-chessboard of size $a \times k$, recalling that $a$ is even
gives us the conclusion that this shape(s) can be tiled due to claim 2.
Here is an illustration of the idea.
\begin{center}
\begin{asy}
unitsize(0.7cm);   // adjust printed size
real s = 1;        // square side in coordinate units

// pens and styles
pen darkSquare  = black;
pen lightSquare = white;
pen borderPen   = black + linewidth(0.6);

pen redFill  = rgb(0.86, 0.20, 0.20) + opacity(0.85);
pen blueFill = rgb(0.15, 0.45, 0.85) + opacity(0.85);
pen circleBorder = black + linewidth(0.5);

real radius = 0.28; // circle radius (in square units)

// helper predicates (i = file 0..7 for a..h, j = rank 0..7 for 1..8)
// removed rectangle: files b..d => i in [1..3], ranks 4..7 => j in [3..6]
bool isRemoved(int i, int j) {
  return (i >= 1 && i <= 3 && j >= 3 && j <= 6);
}
// blue region: files e..h (i=4..7), ranks 4..7 (j=3..6)
bool isBlueRegion(int i, int j) {
  return (i >= 4 && i <= 7 && j >= 3 && j <= 6);
}
// red region: file a (i=0), ranks 4..7 (j=3..6)
bool isRedRegion(int i, int j) {
  return (i == 0 && j >= 3 && j <= 6);
}

// draw the chessboard, skipping removed squares
for (int i = 0; i < 8; ++i) {
  for (int j = 0; j < 8; ++j) {
    if (isRemoved(i,j)) continue;    // omit removed squares
    path square = (i,j) -- (i+1,j) -- (i+1,j+1) -- (i,j+1) -- cycle;
    if ((i + j) % 2 == 0) fill(square, darkSquare);
    else fill(square, lightSquare);
    draw(square, borderPen);
  }
}

// draw colored circles on top: blue in e4..h7, red in a4..a7
for (int i = 0; i < 8; ++i) {
  for (int j = 0; j < 8; ++j) {
    if (isRemoved(i,j)) continue;  // no circle in removed squares
    pair center = (i + 0.5, j + 0.5);
    if (isBlueRegion(i,j)) {
      fill(circle(center, radius), blueFill);
      //draw(circle(center, radius), circleBorder);
    } else if (isRedRegion(i,j)) {
      fill(circle(center, radius), redFill);
      //draw(circle(center, radius), circleBorder);
    }
  }
}

draw((1,3)--(4,3)--(4,7)--(1,7)--cycle, red+linewidth(0.6));

// ---------------------------
// measurement arrow: "<---->" with "2n squares" under the board
// center the arrow under the chessboard; arrow spans 2*n squares
int n = 4;
real arrowY = -0.55;                  // vertical position (below file labels)
real arrowHead = 0.12;                // arrowhead size (in coord units)
real span = 2*n;                      // span in square units
real x0 = (8 - span)/2.0;             // left x (centers it under the 8-wide board)
real x1 = x0 + span;                  // right x

// main line
draw((x0, arrowY) -- (x1, arrowY), borderPen + linewidth(0.8));

// left arrowhead (filled triangle)
fill((x0, arrowY) -- (x0 + arrowHead, arrowY + arrowHead/2) -- (x0 + arrowHead, arrowY - arrowHead/2) -- cycle, borderPen);
draw((x0, arrowY) -- (x0 + arrowHead, arrowY + arrowHead/2), borderPen);
draw((x0, arrowY) -- (x0 + arrowHead, arrowY - arrowHead/2), borderPen);

// right arrowhead
fill((x1, arrowY) -- (x1 - arrowHead, arrowY + arrowHead/2) -- (x1 - arrowHead, arrowY - arrowHead/2) -- cycle, borderPen);
draw((x1, arrowY) -- (x1 - arrowHead, arrowY + arrowHead/2), borderPen);
draw((x1, arrowY) -- (x1 - arrowHead, arrowY - arrowHead/2), borderPen);

// centered label under the arrow
label("$2n$ squares", ((x0 + x1)/2.0, arrowY - 0.34), fontsize(10pt));
\end{asy}
\end{center}

Finally to finish the proof we observe that all there is left to tile, is
squares that lie directly above or below the $a\times b$ hole in the
chessboard. To tile these shape(s) we can just make a rectangle with
all of the squares that lie (above or below) the hole, such a rectangle will
have exactly $2n$ squares on at least one of it's sides and so it is
possible to tile it with dominoes again due to claim 2. $\square$ \\

Now let us return to the original problem. I will describe the
high-level strategy. We will
consider the "sub-chessboard" that is defined by these two removed squares.
We will show that this sub-chessboard with it's two corner squares removed
can be tiled using claim 3, then finally we will show that the remaining
shape "a chessboard with a hole in it" can be tiled using claim 4.
\vspace{0.2cm}

Suppose we have removed two
squares of distinct colours on an $8\times 8$ chessboard,
$(x_1, y_1)$ and $(x_2, y_2)$. Note that the two removed squares,
define a sub-chessboard of dimensions
$(|x_1 - x_2| + 1) \times (|y_1 - y_2| + 1)$.
Due to claim 1 these two squares have
different weights. We show that exactly one of the quantities
$$
|x_1 - x_2| + 1, \quad |y_1 - y_2| + 1
$$
must be even. It suffices to show that their sum is
$\equiv 1 \ (\text{mod } 2)$. We have
\[
\begin{aligned}
&\phantom{\equiv} \left(|x_1 - x_2| + 1\right) +  \left(|y_1 - y_2| + 1\right) \\
&\equiv (x_1 - x_2) + (y_1 - y_2) + 2 \\
&\equiv (x_1 - x_2) + (y_1 - y_2) \\
&\equiv (x_1 + y_1) - (x_2 + y_2) \\
&\equiv 1 \ (\text{mod } 2)
\end{aligned}
\]
where the last line is because the two squares have different weights.
So since this sub-chess board has dimension of the form
$a \times b$, with $a,b \in \mathbb{N}$ and exactly one of $a,b$ even,
The $(|x_1 - x_2| + 1) \times (|y_1 - y_2| + 1)$ chessboard with two
corners removed can be tiled with dominoes by applying claim 3. Finally the
remaining untiled squares form an $8\times 8$ chessboard with a
$(|x_1 - x_2| + 1) \times (|y_1 - y_2| + 1)$ hole in it. Indeed exactly
one of those numbers is even and
$$
\max(|x_1 - x_2| + 1) = \max(|y_1 - y_2| + 1) = 8-1+1=8 \leq 8.
$$
Thus this shape can indeed be tiled due to claim 4. $\blacksquare$
\\

\textit{Comments: I am purposefully hand-waving with all these WLOG's, I'm
not super sure on how to formally justify them tbh. Other than this I
believe the problem is solved now, this one was particularly enjoyable.
I have a feeling that my "chessboard with a hole in it" solution wasn't the
intended one, so im curious to see what other (probably much simpler)
approaches work.}
\end{solution}
\newpage

\begin{problem}{Problem 4}
(a) Suppose that $ABCDEF$ is a regular hexagon. If you connect the
vertices $AC, BD, CE, DF, EA$ and $FB$, then a smaller hexagon is formed.
\vspace{0.2cm}

What is the ratio of the side lengths of the original hexagon to those of
the smaller hexagon?\\

(b) Suppose that $ABCDE$ is a regular pentagon. If you connect vertices
$AC, BD, CE, DA$ and $EB$, then a smaller pentagon is formed. \vspace{0.2cm}

What is the ratio of the side lengths of the original pentagon to those of
the smaller pentagon?
\end{problem}

\begin{solution}{Solution}
(a) Here's a diagram. Let the regular hexagon have the 6-th roots of unity
as vertices.

\begin{center}
\begin{asy}
/* Geogebra to Asymptote conversion, documentation at artofproblemsolving.com/Wiki go to User:Azjps/geogebra */
import graph; size(170);
real labelscalefactor = 0.5; /* changes label-to-point distance */
pen dps = linewidth(0.7) + fontsize(10); defaultpen(dps); /* default pen style */
pen dotstyle = black; /* point style */
real xmin = -2.703120263794179, xmax = 3.7752900936807827, ymin = -2.5178424429372392, ymax = 1.9445066218254417;  /* image dimensions */

pair A = (-0.5,0.8660254037844387), B = (0.5,0.8660254037844386), C = (1.,0.), D = (0.5,-0.8660254037844386), F = (-1.,0.);

filldraw(arc(A,0.2501316740337822,-90.,-30.)--(-0.5,0.8660254037844387)--cycle, red + opacity(0.10000000149011612),  red);
filldraw(arc((0.5,0.28867513459481287),0.2501316740337822,450.,510.00000000000006)--(0.5,0.28867513459481287)--cycle, red + opacity(0.10000000149011612),  red);
/* draw figures */
draw(A--B);
draw(B--C);
draw(C--D);
draw(D--(-0.5,-0.8660254037844384));
draw((-0.5,-0.8660254037844384)--F);
draw(F--A);
draw(A--C);
draw((0.26500790044202716,0.4590071479727455)--(0.2349920995579733,0.4070182558116933));
draw(C--(-0.5,-0.8660254037844384));
draw((0.2650079004420266,-0.4590071479727452)--(0.23499209955797273,-0.40701825581169304));
draw(D--F);
draw(F--B);
draw(B--D);
draw(A--(-0.5,-0.8660254037844384));
draw((-0.4699841991159463,0.)--(-0.5300158008840541,0.));
/* dots and labels */
dot(C,dotstyle);
dot(B,dotstyle);
dot(A,dotstyle);
dot(F,dotstyle);
dot((-0.5,-0.8660254037844384),dotstyle);
dot(D,dotstyle);
dot(A,dotstyle);
label("$A$", (-0.48195099837419186,0.9139641248062578), NE * labelscalefactor);
dot(B,dotstyle);
label("$B$", (0.518575697760937,0.9139641248062578), NE * labelscalefactor);
dot(C,dotstyle);
label("$C$", (1.0188390458285015,0.04850853264937002), NE * labelscalefactor);
dot(D,dotstyle);
label("$D$", (0.518575697760937,-0.8169470595075177), NE * labelscalefactor);
dot((-0.5,-0.8660254037844384),dotstyle);
label("$E$", (-0.48195099837419186,-0.8169470595075177), NE * labelscalefactor);
dot(F,dotstyle);
label("$F$", (-0.9822143464417563,0.04850853264937002), NE * labelscalefactor);
label("$60^\circ$", (-0.39190359572203026,0.5637797811589621), NE * labelscalefactor,red);
label("$60^\circ$", (0.2884545576498574,0.5187560798328813), NE * labelscalefactor,red);
clip((xmin,ymin)--(xmin,ymax)--(xmax,ymax)--(xmax,ymin)--cycle);
/* end of picture */
\end{asy}
\end{center}

Since this is a regular hexagon, rotating the diagram by $120^\circ$
anti-clockwise about the origin will send $C \rightarrow A$ and
$A \rightarrow E$. Similarly rotating the diagram by $120^\circ$
clockwise about the origin will send $E \rightarrow A$ and
$C \rightarrow E$. Since the length of a segment is unchanged
under rotation this allows us to conclude that $AC = CE = AE$.
\vspace{0.2cm}

Now since $\triangle ACE$ is equilateral, we have
$\angle CAE = 60^\circ$. Then note that $AE \parallel BD$, because they map
to each other under a reflection in the line $x=0$. Thus by alternate angles
we deduce that $\angle AHB = 60^\circ$ where $H = AC \cap BD$.
Now since the hexagon $ABCDEF$ has side length 1, we deduce that
$$
AH = \frac{1}{\sin 60^\circ} = \frac{2}{\sqrt3}.
$$
Now where $X = AC \cap BF$, we note that $X$ lies on the midpoint of
$AH$, because $BF$ is just the reflection of $AC$ over $x=0$ and so the only
point at which they intersect would be at $x=0$. Thus
$$
HX = \frac12 \frac{2}{\sqrt3} = \frac{1}{\sqrt3}
$$
and so the desired ratio is $1 : \frac{1}{\sqrt3} = \sqrt3 : 1$. $\square$
\\

\newpage
(b) Here's a diagram.

\begin{center}
\begin{asy}
/* Geogebra to Asymptote conversion, documentation at artofproblemsolving.com/Wiki go to User:Azjps/geogebra */
import graph; size(170);
real labelscalefactor = 0.5; /* changes label-to-point distance */
pen dps = linewidth(0.7) + fontsize(10); defaultpen(dps); /* default pen style */
pen dotstyle = black; /* point style */
real xmin = -2.160089228990391, xmax = 1.7239076975923278, ymin = -1.1511653384781935, ymax = 1.51385409114626;  /* image dimensions */

pair B = (0.,1.), C = (0.9510565162951536,0.30901699437494723), D = (0.587785252292473,-0.8090169943749476), A = (-0.9510565162951535,0.30901699437494745), G = (-0.22451398828979274,0.30901699437494734), H = (0.2245139882897927,0.30901699437494723);

filldraw(arc(B,0.19883076574159989,-144.,-108.)--(0.,1.)--cycle, red + opacity(0.10000000149011612),  red);
filldraw(arc(B,0.19883076574159989,-108.,-72.)--(0.,1.)--cycle, red + opacity(0.10000000149011612),  red);
/* draw figures */
draw(A--B);
draw(B--C);
draw(C--D);
draw(D--(-0.5877852522924732,-0.8090169943749473));
draw((-0.5877852522924732,-0.8090169943749473)--A);
draw(A--C);
draw(B--D);
draw(C--(-0.5877852522924732,-0.8090169943749473));
draw(D--A);
draw((-0.5877852522924732,-0.8090169943749473)--B);
/* dots and labels */
dot(B,dotstyle);
dot(A,dotstyle);
dot((-0.5877852522924732,-0.8090169943749473),dotstyle);
dot(D,dotstyle);
dot(C,dotstyle);
dot(B,dotstyle);
label("$B$", (0.01485941919762576,1.0349572700986107), NE * labelscalefactor);
dot(C,dotstyle);
label("$C$", (0.9656319899256397,0.3444722472505079), NE * labelscalefactor);
dot(D,dotstyle);
label("$D$", (0.600506401927429,-0.7725951457341191), NE * labelscalefactor);
dot((-0.5877852522924732,-0.8090169943749473),dotstyle);
label("$E$", (-0.574402668363843,-0.7725951457341191), NE * labelscalefactor);
dot(A,dotstyle);
label("$A$", (-0.9359131515303882,0.3444722472505079), NE * labelscalefactor);
label("$36^\circ$", (-0.20927708036563228,0.7602093028920358), NE * labelscalefactor,red);
label("$36^\circ$", (-0.042982258109021475,0.7095978352487193), NE * labelscalefactor,red);
dot(G, dotstyle);
label("$G$", (-0.20927708036563228,0.33724203758717697), NE * labelscalefactor);
dot(H, dotstyle);
label("$H$", (0.2389959187608838,0.33724203758717697), NE * labelscalefactor);
clip((xmin,ymin)--(xmin,ymax)--(xmax,ymax)--(xmax,ymin)--cycle);
/* end of picture */
\end{asy}
\end{center}

Firstly because $\angle EAB = 108^\circ$ and $AB = AE$, $\triangle ABE$ is
isoceles and so $\angle ABE = 36^\circ$. By a similar arguement
$\angle DBC = 36^\circ$ which forces $\angle GBH = 36^\circ$ since a
regular pentagon have interior angles of $108^\circ$. \vspace{0.2cm}

Now if we let the regular pentagon have side length $1$, we can use sine
rule to find that
$$
GB = \frac{\sin 36^\circ}{\sin 108^\circ} =
\frac{\sin 36^\circ}{\sin (180-108)^\circ} =
\frac{\sin 36^\circ}{\sin 72^\circ} =
\frac{\sin 36^\circ}{\cos 36^\circ} = \tan 36^\circ .
$$
Then again by sine rule
$$
GH = \frac{\sin 36^\circ}{\sin 72^\circ} GB =
\frac{\sin 36^\circ}{\cos 36^\circ} GB = \tan^2 36^\circ.
$$
So the desired ratio is $1 : \tan^2 36^\circ$. \vspace{0.2cm}

Now we attempt to evaluate $\tan^2 36^\circ$. We can use De Moivre's
theorem as follows. Let $\theta = \pi/5$.
\[
\begin{aligned}
\sin 5\theta &= \text{Im } (e^{i\theta})^5\\
&= \text{Im } \sum_{r=0}^5 (i\sin \theta)^r \cos^{5-r} \theta \\
&= x(16x^4 - 20x^2 + 5) \\
\end{aligned}
\]
where $x = \sin \theta$. The roots of this polynomial correspond to
$$
x = \sin \frac{k \pi}{5}, \quad (k = 0, 1, 2, 3, 4).
$$
Clearly $k=0$ corresponds to the lone factor of $x$ so let's ignore it.
If we solve that quartic in terms of $x^2$ we get
$$
x^2 = \frac{5 \pm \sqrt{5}}{8}.
$$
Since $\sin \theta$ is positive on the interval $(0, \pi)$, the 4 roots must
be
$$
\sqrt{\frac{5-\sqrt{5}}{8}}, \ \sqrt{\frac{5-\sqrt{5}}{8}}, \
\sqrt{\frac{5+\sqrt{5}}{8}}, \ \sqrt{\frac{5+\sqrt{5}}{8}}.
$$
Now if you look at the graph of $\sin x$ you can see that $\sin \pi/5$ and
$\sin 4\pi/5$
should be amongst the smallest of these values so
$$
\sin \frac{\pi}{5} = \sqrt{\frac{5-\sqrt5}{8}}.
$$
Now since both $\sin \pi/5$ and $\cos \pi/5$ are positive we can compute
$$
\cos\frac{\pi}{5} = \sqrt{1 - \left(\sqrt{\frac{5-\sqrt5}{8}}\right)^2} =
\sqrt{\frac{3+\sqrt{5}}{8}}.
$$
Now recalling that $\tan^2 \theta = \sec^2 \theta -1$ we have
$$
\tan^2 \frac{\pi}{5} = \frac{8}{3+\sqrt5} - 1 = 5 - 2\sqrt5
$$
and thus the desired ratio turns out to be $1 : (5-2\sqrt5)$. $\blacksquare$
\end{solution}
\newpage

\begin{problem}{Problem 5}
(a) Show that
$$
\binom{n}{r} = \binom{n}{n-r}.
$$

(b) Show that
$$
\binom{n}{0}^2 + \binom{n}{1}^2 + \cdots + \binom{n}{n}^2 = \binom{2n}{n}.
$$
\end{problem}

\begin{solution}{Solution}
(a) We note that
$$
\binom{n}{n-r} = \frac{n!}{(n-r)!(n - (n-r))!} = \frac{n!}{(n-r)!r!} =
\binom{n}{r}. \ \square
$$

(b) This result is a direct application of Vandermonde's identity.
\vspace{0.2cm}

\textit{Theorem. Vandermonde's Identity.} Let $q \leq \min\{m, n\}$ be
positive integers. Then
$$
\sum_{r=0}^{q} \binom{m}{r} \binom{n}{q-r} = \binom{m+n}{q}.
$$

\textit{Proof:} Consider
\[
\begin{aligned}
(1+x)^m (1+x)^n &= (1+x)^{m+n} \\
\left[ \sum_{r=0}^m \binom{m}{r} x^r \right]
\left[ \sum_{k=0}^n \binom{n}{k} x^k \right] &=
\sum_{t=0}^{m+n}\binom{m+n}{t} x^t\\
\end{aligned}
\]
Now compare the coefficients of $x^q$,
$$
\sum_{r=0}^q \binom{m}{r} \binom{n}{q-r} = \binom{m + n}{q}. \ \square
$$

If we apply this theorem with $(q,m,n) = (n, n, n)$ we get
$$
\sum_{r=0}^n \binom{n}{r}\binom{n}{n-r} = \binom{2n}{n}.
$$
Now using the result from part (a) we can re-write this as
$$
\sum_{r=0}^n \binom{n}{r}^2 = \binom{2n}{n}. \ \blacksquare
$$
\end{solution}
\newpage

\begin{problem}{Problem 6}
2021 has the property that it is formed by concatenating two consecutive
numbers (20 and 21), and it is also the product of two numbers that differ
by 4 ($43 \times 47$).
\vspace{0.2cm}

Find all other 4 digit numbers with these properties.
\end{problem}

\begin{solution}{Solution}
I assume that when we say concatenating two consecutive numbers we write
$\overline{x(x+1)}$. If writing $\overline{(x+1)x}$ is allowed the strategy
detailed below doesn't seem to work. \vspace{0.2cm}

Call a number \textit{special} if it satisfies these properties.
Indeed the number that is formed by concatenating two consecutive numbers
$\overline{x(x+1)}$ can be expressed as $100x + (x+1) = 101x + 1$.
The problem reduces to finding solutions to
$$
101x + 1 = k(k+4)
$$
for $k,x \in \mathbb{N}$.
We can complete the square to find that
$$
(k+2)^2 - 5 = 101x.
$$
Subtract $20\cdot 101$ from both sides, (see below for the motivation)
$$
(k+2)^2 - 5 - 20\cdot 101 = (k+2)^2 - 2025 = 101(x-20)
$$
Noting that $2025 = 45^2$ we complete the square,
$$
(k-43)(k+47) = 101(x-20)
$$
Since $k(k+4)$ must be a 4 digit number we require $k \leq 99$.
Now we note that $101$ must divide one of $(k+47)(k-43)$.

\begin{itemize}
\item If $101 \mid (k+47)$ then using $1 \leq k \leq 99$ it's easy to show that only
$k + 47 = 101 \implies k = 54$.
\item If $101 \mid (k-43)$ then using $1\leq k \leq 99$ it's easy to show that only
$k-43 = 0 \implies k = 43$.
\end{itemize}

From here it can be checked that $(k,x) = (54, 31), (43, 20)$ are the
only solutions. So the only 4 digit special numbers are $3132, \ 2021$
(assuming the $\overline{x(x+1)}$ convention.) $\blacksquare$
\vspace{1.0cm}

\textit{Motivation:} Once I reached the equation $(k+2)^2 - 5 = 101x$, I
spent a long time trying to find solutions to
$(k+2)^2 \equiv 5 \ (\text{mod } 101)$, but I wasn't able to find anything
here. Then after thinking a while about how I can use the fact that $101$
is prime, I thought if I can factor the LHS I might be able to force
unique values of $k$ since the range for $1 \leq k \leq 99$ is fairly
small. I looked for $a\in \mathbb{N}$ such that $5 + 101a$ is square
so that I could factor the LHS with difference of two squares as desired.
I wrote a quick python script and found the first two solutions $a=20,31$.
\vspace{0.2cm}

Re-reading over this I realise that finding $a$ such that $101a + 5$ is
square is essentially the same as my original problem of finding
solutions to $(k+2)^2 \equiv 5 \ (\text{mod } 101)$. Right now I can't see
any better method other than trial and error to solve this congruence, or
guessing to solve here.
\end{solution}
\newpage

\begin{problem}{Problem 7}
Suppose that $p_1, p_2, \ldots, p_n$ are all primes.\vspace{0.2cm}

(a) Consider the quantity $p_1 p_2 \cdots p_n + 1$. Could this be
divisible by any of the $p_i$ ? \vspace{0.2cm}

(b) Using this result prove that there are infinitely many primes.
\vspace{0.2cm}

(c) By considering the quantity $4p_1 p_2 \cdots p_n - 1$, show that there
are infinitely many primes of the form $4k+3$.
\end{problem}

\begin{solution}{Solution}
(a) Notice that $p_1 p_2 \cdots p_n + 1 \equiv 1 \ (\text{mod } p_i)$.
Using the Euclidean algorithm,
$$
\gcd(p_1 p_2 \cdots p_n + 1, p_i) =
\gcd(1, p_i) = 1.
$$
So $p_i$ never divides $p_1 p_2 \cdots p_n + 1$. $\square$
\\

(b) If there are finitely many primes $p_1, p_2, \ldots, p_n$ and
$P = p_1 p_2 \cdots p_n + 1$ is coprime to all of the primes $p_i$, then
due to the fundamental theorem of arithmetic $P$ must have some prime
factors, in particular it's prime factors aren't among
$p_1, \ldots, p_n$. But this contradicts the assumption that
$p_1,\ldots, p_n$ are all the primes. $\square$ \\

(c) Suppose there are finitely many primes congruent to
$3 \ (\text{mod } 4)$, say $q_1, q_2, \ldots q_n$. Consider the quantity,
$$
P = 4q_1 q_2 \cdots q_n - 1.
$$
Note that for every $q_i$ we have $P \equiv -1 \ (\text{mod } q_i)$, and
so by the Euclidean algorithm,
$$
\gcd(4q_1 q_2 \cdots q_n - 1, q_i) = \gcd(-1, q_i) = 1.
$$
So for every $q_i$, we have $q_i \nmid P$. By the fundamental theorem of
arithmetic $P$ factors into a set of primes. Since $P$ is odd, all of these
primes are congruent to $\pm 1 \ (\text{mod } 4)$. However since $P$ itself
is congruent to $-1 \ (\text{mod } 4)$, $P$ must have
$2k+1, (k \in \mathbb{Z}_{\geq0})$
pairwise distinct prime factors that are congruent to
$-1 \ (\text{mod } 4)$. Moreover since $q_i \nmid P$, this means that there
exists at least one new prime $q_{n+1} \equiv -1 \ (\text{mod } 4)$ which
is not a part of the collection $q_1, \ldots, q_n$, a contradiction.
$\square$\\

\textit{Comments:} I first saw the problem in part (c) at the Winchester
maths summer camp in Year 12. At that time I had just started discovering
contest math, so when the teacher (Dominic Rowland) tried explaining this
proof to me I was lost and couldn't understand a thing. Good to know that
now I can actually understand and explain the proof. \\

Here's a note explaining the Euclidean algorithm. Basically
suppose we have $a,b \in \mathbb{Z}$ satisfying
$a \equiv r \ (\text{mod } b)$. Then
$$
\gcd(a,b) = \gcd(r, b).
$$
This is because the gcd, $g$ must be a divisor of $b$. Since we also have
$g \mid a$, we require
$$
\frac{a}{g} = \frac{kb + r}{g} = k\left(\frac{b}{g}\right) + \frac{r}{g}
$$
to be an integer, which happens iff $g \mid r$.
\end{solution}

\end{document}
